% diagrams/firstwords/lucia-nada.tex -- Lucía nada en el río. Para
% "completa": falta la roca grande en medio del río, donde se sienta al
% sol, que se dibuja (el hueco de la derecha).
\begin{tikzpicture}[dibujofw]
  % el agua
  \filldraw[fill=white] (-6.5,-2.5) -- (-6.5,1.0) -- plot[domain=-6.5:6.5, samples=60, smooth] (\x,{1.0+0.14*sin(\x*140)}) -- (6.5,-2.5) -- cycle;
  % los juncos de la orilla
  \foreach \x/\h in {-6.2/3.0, -5.8/3.6, -5.4/2.8} {\draw (\x,1.0) -- (\x,\h);
    \filldraw[fill=white] (\x,{\h+0.3}) ellipse (0.14 and 0.4);}
  % Lucía, nadando: la cabeza, un brazo fuera del agua y las ondas
  \filldraw[fill=white, rounded corners=2mm] (-1.3,1.2) -- (-2.8,2.9) -- (-2.5,3.2) -- (-0.9,1.6) -- cycle;
  \filldraw[fill=white] (-2.75,3.15) circle (0.28);
  \begin{scope}[shift={(-0.6,1.9)}, scale=0.95] \cabezaLucia \end{scope}
  \fill[white] (-2.1,0.55) rectangle (0.9,0.95);
  \draw (-2.3,0.95) .. controls (-1.5,0.7) and (0.3,0.7) .. (1.1,0.95);
  \foreach \x in {-3.6,1.8} {\draw (\x,0.55) .. controls ({\x+0.4},0.75) .. ({\x+0.8},0.55);}
  \path (-6.5,0) (6.5,5.0);
\end{tikzpicture}
